\section{Specify precedence in the grammar}\label{expression-grammar}
The grammar uses four named rules. An expression contains terms separated by plus
or minus. A term contains unary operands separated by multiplication or division.
A unary operand is either a sign followed by another unary operand, or a primary.
A primary is a decimal literal or a parenthesized expression.

\begin{table}[htbp]
  \centering
  \caption{Grammar for this exercise. Braces mean zero or more repetitions; a vertical bar separates alternatives. Whitespace may appear between tokens.}
  \tagpdfsetup{table/header-rows={1}}
  \begin{tabular}{ll}
    \toprule
    Rule       & Production                                           \\
    \midrule
    Expression & \texttt{term \{ (\char`+ | \char`-) term \}}         \\
    Term       & \texttt{unary \{ (\char`* | /) unary \}}             \\
    Unary      & \texttt{(\char`+ | \char`-) unary | primary}         \\
    Primary    & \texttt{digits | ( expression )}                     \\
    Digits     & One or more bytes from \texttt{0} through \texttt{9} \\
    \bottomrule
  \end{tabular}
\end{table}
The grammar gives multiplication and division higher precedence than binary plus
and minus: a term finishes its multiplicative work before returning a value to the
expression routine. A unary sign binds to its operand before that operand is used
by a binary operator. Parentheses let a primary contain a complete expression,
regardless of the precedence level outside the group.

The repetition in the expression and term rules is evaluated left to right.
For \texttt{1-2-3}, start with one, subtract two to obtain minus one, and then
subtract three to obtain minus four. For \texttt{8/3/2}, the first quotient is
two and the second is one. Right association would produce a different computation.
Unary signs recurse instead: in \texttt{--5}, the inner minus produces minus five
and the outer minus produces five. Repeated signs do not form a separate decrement
or increment token in this language; \texttt{1++2} means one plus unary-positive two.
